\documentclass[../../project_master.tex]{subfiles}
\begin{document}
\section{Silicon Defect Physics}
Radiation-induced defects in silicon are important because they create electrically active states that modify carrier statistics and transport. In a reduced semiconductor model, the most important consequences are enhanced recombination, enhanced trapping, and increased generation current.

\subsection{Trap and Recombination Centers}
A defect level in the bandgap can capture an electron or a hole and later emit it. When a defect efficiently captures both carrier types, it can serve as a Shockley--Read--Hall recombination center. This reduces the minority-carrier lifetime and is especially important for devices that depend on long carrier lifetimes or efficient charge collection.

\subsection{Lifetime Degradation}
If the defect population grows with radiation exposure, the effective lifetime is expected to fall. In a simple reduced model, the inverse lifetime grows approximately linearly with defect density, which motivates the working expression
\begin{equation}
\frac{1}{\teff} = \frac{1}{\tzero} + K_\tau \Nd.
\end{equation}

\subsection{Transport Consequences}
Once the lifetime decreases, the diffusion length also decreases through
\begin{equation}
L = \sqrt{D\,\teff}.
\end{equation}
The resulting reduction in diffusion length is one of the cleanest ways to connect microscopic defect physics to device-level signal degradation.
\end{document}
